% kepler_newton_ML.tex
% "Kepler and Newton as Machine Learners:
%  From Tycho Brahe's Data to Universal Gravitation"
% Thomas J. Sargent, 2026

\documentclass[12pt]{article}

\usepackage{amsmath,amssymb}
\usepackage{booktabs}
\usepackage{array}
\usepackage{multirow}
\usepackage[margin=1.2in]{geometry}
\usepackage{setspace}
\usepackage{microtype}
\usepackage{hyperref}
\usepackage{natbib}
\usepackage{parskip}

\DeclareMathOperator*{\argmin}{arg\,min}

\onehalfspacing

\hypersetup{
  colorlinks  = true,
  linkcolor   = blue,
  citecolor   = blue,
  urlcolor    = blue
}

\title{Kepler and Newton as Machine Learners:\\
       From Tycho Brahe's Data to Universal Gravitation}
\author{Thomas J.\ Sargent}
\date{March 2026}

\begin{document}
\maketitle

\begin{abstract}
Between 1600 and 1609, Johannes Kepler performed history's first large-scale
nonlinear least-squares regression: fitting ellipse parameters to Tycho Brahe's
Mars observations, which achieved then-unprecedented $2'$-arc\-minute precision.
Eighty years later, Isaac Newton derived the law of universal gravitation from
Kepler's three empirical laws by a sequence of deductions that constitute a
hypothetico-deductive inference chain.  This paper examines the two-century
episode from the perspective of modern statistics and machine learning.
Kepler's problem was nonlinear parameter estimation: the elliptical model
$r(\theta)=a(1-e^{2})/(1+e\cos\theta)$ was new, the parameters $(a,e)$
unknown, and the loss function the sum of squared angular residuals.
The optimal Mars parameters~--- $a=1.524$\,AU, $e=0.0934$~---
can be recovered in seconds by nonlinear least squares or MCMC; what required
Kepler nearly a decade was the prior discovery of the correct functional form.
The 8-arc\-minute residual of the circular hypothesis played the role of a
four-sigma outlier: statistically inconsistent with Brahe's measurement
precision, it forced a model expansion from circular to elliptical orbits.
Kepler's Third Law $T^{2}\propto a^{3}$, which he discovered inductively across
the five planets then known, is a log-log OLS regression: the slope estimate
is $\hat\beta=1.500$ and $R^{2}=0.99998$.
Newton's contribution was an inverse problem: given Kepler's three empirical
regularities as constraints, identify the unique force law $F(r)$ that
produces them.  The answer, $F\propto r^{-2}$, follows from elementary
calculus and the Binet equation; the Moon test then confirmed out-of-sample
generalization from planetary to terrestrial gravitation.
Modern AI systems~---
\textit{AI~Feynman} \citep{Udrescu2020},
\textit{AI-Newton} \citep{FangEtAl2023},
and the deep-learning network of \citet{Lemos2023}~---
reconstruct the same laws from trajectory data with no prior physics knowledge.
The contrast illuminates both the power of theory-informed inference and the
irreducible role of conceptual insight in scientific discovery.
\end{abstract}

\tableofcontents

%======================================================================
\section{Introduction}

A modern data scientist confronted with a table of (orbital period, semi-major
axis) pairs would recognise Kepler's Third Law as a two-variable power-law
regression, estimated in about one line of Python.
A data scientist given a time series of planetary positions would frame the
orbit-fitting problem as nonlinear least squares with two free parameters.
Neither would need Kepler's conceptual leap~---the abandonment of circles in
favour of ellipses~---because they would already know the functional form.
What Kepler \emph{did} need, and what took him nine years and hundreds of pages
of \emph{Astronomia Nova} \citep{Donahue1992}, was to discover the model class
before estimating its parameters.

The episode repays examination from a statistical perspective for three reasons.

First, the \emph{data}.  Tycho Brahe's naked-eye observations of Mars,
accumulated over twenty years at Uraniborg, achieved $2'$\,(arc\-minute)
precision~---roughly ten times better than anything available to Copernicus~---
and were thus sufficient to distinguish between competing orbit models.
Brahe's archive was the training set; its precision determined the statistical
power available to any estimation procedure.

Second, the \emph{model-selection criterion}.
Kepler's pivotal insight was not a computation but a residual diagnostic:
when the Vicarious Hypothesis (an eccentric circle fitted to Brahe's data)
predicted Mars longitudes within $2'$ but left an irreducible $8'$ error at
certain orbital phases, Kepler recognised that $8'=4\sigma$ was statistically
incompatible with measurement noise.
He could not ignore it, and so was forced to enlarge the model class.
The ellipse~---the next simplest closed curve after the circle~---fitted the
residuals to within the noise floor, and no further expansion was needed.
This is textbook model selection: reject the null model when the residual
exceeds a critical threshold, expand the hypothesis class, and stop as soon
as the residuals are consistent with noise.

Third, the \emph{identification problem}.
Newton's contribution was not more data or a better fitting algorithm.
It was the recognition that Kepler's three empirical laws, taken together,
over-identify a unique force law.
The area law (equal areas in equal times) implied a central force.
The harmonic law ($T^{2}\propto a^{3}$) implied an inverse-square dependence
on distance.
The fact that orbits are ellipses~--- not parabolas or other conics~---
confirmed the inverse-square law by the converse of Proposition~11, Book~I
of the \emph{Principia}.
The Moon test extended the law from the solar system to the Earth–Moon system,
establishing universality across seven orders of magnitude in distance.

Seen through a machine-learning lens, Brahe was the data engineer,
Kepler was the model builder who performed nonlinear estimation under a
parsimony constraint, and Newton solved the inverse problem of recovering the
latent physical law from Kepler's compressed representation of the data.
Modern AI systems have now automated all three steps.
The remainder of this paper develops each step in detail, provides numerical
worked examples, and closes with a comparison of Kepler and Newton's
reasoning with the methods of contemporary machine learning.

%======================================================================
\section{Tycho Brahe's Data and the Mars Problem}

\subsection{Observational Precision}

Tycho Brahe (1546--1601) constructed the most precise pre-telescopic
observational instruments in history at his observatory Uraniborg on the
island of Hven.  His quadrants and sextants measured celestial angles to
approximately $2'$ (arc\-minutes), an order of magnitude more precisely than
his predecessors.  Over two decades he accumulated positions for all of the
then-known planets, the Moon, and over a thousand fixed stars.

Two arc\-minutes is a small angle~---$(2/60)\times(\pi/180)=5.8\times10^{-4}$
radians~---but it is large enough to matter.
At Mars's mean distance of $1.524$\,AU from the Sun, a $2'$ angular error
corresponds to a linear uncertainty of
\[
   \delta r = 1.524 \times 5.8\times10^{-4}\,\text{AU}
            = 8.8\times10^{-4}\,\text{AU} \approx 0.001\,\text{AU},
\]
or about $130{,}000$\,km.  This is the irreducible noise floor in Kepler's
estimation problem.

\subsection{The Method of Oppositions}

Kepler used Mars observations taken at \emph{opposition}~---
moments when the Sun, Earth, and Mars are collinear, with Earth in
the middle.  At opposition, the geocentric and heliocentric longitudes
of Mars coincide (because the Earth--Sun--Mars angle is zero), so no
geometric correction for the Earth's position is needed.
Kepler identified twelve usable oppositions in Brahe's archive spanning 1580
to 1600.  Each opposition supplied one data point $(r_{i},\theta_{i})$:
the heliocentric distance $r_{i}$ and the heliocentric true anomaly
$\theta_{i}$ of Mars.

\subsection{Table~1: Selected Brahe--Kepler Opposition Data for Mars}

Table~\ref{tab:mars_data} lists five of Kepler's Mars opposition data
points in the form convenient for parameter estimation.
The values are approximate but historically representative, with added noise
consistent with the $2'$ precision of Brahe's instruments.

\begin{table}[h]
\centering
\renewcommand{\arraystretch}{1.25}
\caption{Selected Mars opposition heliocentric coordinates derived from
  Brahe's observations.  Distances in AU; true anomaly $\theta$ measured
  from perihelion.  Noise is of order $0.001$--$0.002$\,AU, consistent with
  Brahe's $2'$ angular precision at $1.5$\,AU distance.}
\label{tab:mars_data}
\begin{tabular}{@{}ccc@{}}
\toprule
$i$ & $\theta_{i}$ (degrees) & $r_{i}$ (AU) \\
\midrule
1 &   0  & 1.382 \\
2 &  75  & 1.477 \\
3 & 150  & 1.641 \\
4 & 225  & 1.619 \\
5 & 300  & 1.442 \\
\bottomrule
\end{tabular}
\end{table}

%======================================================================
\section{Kepler's Parametric Model}

\subsection{The Vicarious Hypothesis and Its Failure}

Kepler's initial working model for Mars was the \emph{Vicarious Hypothesis}:
an eccentric circle, with an equant point offset from the centre of the
circle.  This model had been the standard since Ptolemy.
Fitted to oppositions, it predicted heliocentric longitudes within $2'$~---
inside Brahe's error budget~---and was thus statistically acceptable on
longitude residuals alone.

The problem appeared when Kepler checked the \emph{latitudes} and distances.
At orbital phases halfway between opposition and conjunction, the circular
model produced errors of $8'$~---four times Brahe's precision.
Kepler wrote in \emph{Astronomia Nova} (Chapter 19):
\begin{quote}
  \emph{``If thou, dear reader, art bored with this wearisome method of
  calculation, take pity on me who had to go through with at least seventy
  repetitions of it\ldots{} But if you are wearied by this prolixity,
  you will be still more wearied when I tell you that this hypothesis is
  nevertheless wide of the truth.''}
\end{quote}
And, crucially (Chapter 19):
\begin{quote}
  \emph{``These eight minutes alone will have led the way to the reformation
  of all of astronomy.''}
\end{quote}

\subsection{The Ellipse}

After confirming that no circular hypothesis could explain Brahe's data,
Kepler tried the oval and ultimately the ellipse.
The ellipse, in polar coordinates centred at one focus (the Sun), is
\begin{equation}\label{eq:ellipse}
  \boxed{r(\theta) = \frac{a(1-e^{2})}{1+e\cos\theta}}
\end{equation}
where $a$ is the semi-major axis, $e\in(0,1)$ is the eccentricity,
and $\theta$ is the true anomaly measured from perihelion.
Two free parameters, $(a,e)$, specify the size and shape of the orbit.
A third parameter, the argument of perihelion, specifies the orientation
of the ellipse in space, but this can be absorbed into the definition
of $\theta=0$.

\subsection{Kepler's Equation}

The ellipse specifies the orbit's \emph{shape}: where the planet is at a
given true anomaly.  To predict the planet's position at a given \emph{time},
one needs to connect $\theta$ to time via the eccentric anomaly $E$ and the
mean anomaly $M$.  The key relation is \emph{Kepler's equation}:
\begin{equation}\label{eq:kepler}
  \boxed{M = E - e\sin E}
\end{equation}
where $M = 2\pi(t-t_{0})/T$ increases linearly with time and $E$ is
related to $\theta$ by $\tan(\theta/2)=\sqrt{(1+e)/(1-e)}\,\tan(E/2)$.
Kepler's equation~\eqref{eq:kepler} is transcendental in $E$ and must be
solved iteratively, an additional nonlinearity that Kepler handled by
successive approximation~---a procedure equivalent to Newton's method applied
to $g(E)=M-E+e\sin E=0$.

%======================================================================
\section{The Eight-Minute Pivot: Model Selection by Residual Diagnostics}
\label{sec:8min}

\subsection{The Circular Model as a Null Hypothesis}

The Vicarious Hypothesis~---optimally fitted to Brahe's data~---predicted
Mars longitudes within $2'$ at opposition epochs, consistent with measurement
noise.  The hypothesis could not be rejected on longitude residuals.

The problem was that it predicted Mars distances incorrectly.
Converting distance errors to angular terms at Mars's $1.524$\,AU distance,
the deviation reached $8'$.  Since Brahe's measurement precision was $\sigma
\approx 2'$, the circular model's maximum residual was
\[
   t_{\max} = \frac{8'}{2'} = 4\,\text{ standard deviations}.
\]
Under a Gaussian noise model, a four-sigma deviation occurs with probability
$p = 2\Phi(-4) \approx 6.3\times10^{-5}$.  The circular hypothesis was
decisively rejected.

\subsection{Model Expansion}

The rejection of the circular null drove Kepler to enlarge his model class in
the direction of greater flexibility.
The sequence of candidates, ordered by complexity, was:
\begin{enumerate}
  \item Circle (one shape parameter: radius $a$).
  \item Eccentric circle with equant (two geometric parameters: eccentricity
        $e$, equant offset $d$).  This is the Vicarious Hypothesis.
  \item Oval (one additional shape parameter).
  \item Ellipse (two parameters: $a$, $e$).
\end{enumerate}
Kepler tried the oval and found it over-corrected at some phases while
under-correcting at others.  The ellipse, with its elegant two-parameter
form, reduced all residuals to within $2'$~---inside the noise floor~---
and passed the residual test.  No richer model was needed.

This is the logic of the \emph{Minimum Description Length} (MDL) principle
\citep{Udrescu2020,SchmidtLipson2009}: prefer the model that minimises the
sum of (i)~the cost of encoding the model's parameters and (ii)~the cost of
encoding the residuals under the model.
The ellipse minimises this sum because it captures the systematic structure
in Brahe's data with only two parameters, leaving only random noise in the
residuals.

%======================================================================
\section{Nonlinear Least Squares for the Ellipse Parameters}
\label{sec:nls}

\subsection{The NLS Objective}

Given $n$ opposition observations $\{(\theta_{i},r_{i})\}_{i=1}^{n}$ and
an assumed noise model $r_{i}=f(\theta_{i};a,e)+\varepsilon_{i}$ where
$\varepsilon_{i}\overset{iid}{\sim}(0,\sigma^{2})$, the NLS estimator
minimises the sum of squared residuals:
\begin{equation}\label{eq:nls}
  (\hat{a},\hat{e}) = \argmin_{a,e}
    \sum_{i=1}^{n}\bigl[r_{i} - f(\theta_{i};a,e)\bigr]^{2}
\end{equation}
where $f(\theta;a,e)=a(1-e^{2})/(1+e\cos\theta)$, as in
equation~\eqref{eq:ellipse}.

Unlike OLS, the NLS problem~\eqref{eq:nls} is nonlinear in the parameters.
No closed-form solution exists.
The standard algorithm is \emph{Gauss--Newton}: starting from an initial
guess $(a^{(0)},e^{(0)})$, each iteration solves a linearised least-squares
problem.
Let $\mathbf{J}^{(k)}$ be the $n\times2$ Jacobian matrix with
\[
  J_{i1}^{(k)} = \frac{\partial f}{\partial a}\Big|_{\theta_{i},\,a^{(k)},e^{(k)}}
               = \frac{1-e^{2}}{1+e\cos\theta_{i}},
  \qquad
  J_{i2}^{(k)} = \frac{\partial f}{\partial e}\Big|_{\theta_{i},\,a^{(k)},e^{(k)}}
               = \frac{-a(2e+(1+e^{2})\cos\theta_{i})}{(1+e\cos\theta_{i})^{2}}.
\]
The Gauss--Newton update is
\begin{equation}\label{eq:gaussnewton}
  \begin{pmatrix}a^{(k+1)}\\e^{(k+1)}\end{pmatrix}
  = \begin{pmatrix}a^{(k)}\\e^{(k)}\end{pmatrix}
  + \bigl((\mathbf{J}^{(k)})^{\top}\mathbf{J}^{(k)}\bigr)^{-1}
    (\mathbf{J}^{(k)})^{\top}\mathbf{r}^{(k)},
\end{equation}
where $\mathbf{r}^{(k)}$ is the residual vector at the $k$-th iterate.
The algorithm converges in a few iterations when started near the solution.
Kepler's iterative procedure~---adjusting $a$ and $e$ by hand, checking
residuals, and revising~---is logically equivalent to early iterations of
Gauss--Newton, with human judgement substituting for the matrix product.

\subsection{NLS Residual Table for Mars}

Table~\ref{tab:nls_residuals} shows the NLS fit evaluated at the converged
estimates $\hat{a}=1.524$\,AU and $\hat{e}=0.0934$, applied to the five
observations in Table~\ref{tab:mars_data}.

\begin{table}[h]
\centering
\renewcommand{\arraystretch}{1.25}
\caption{NLS residuals for the five Mars opposition observations.
  Fitted model: $\hat{r}(\theta)=\hat{a}(1-\hat{e}^{2})/(1+\hat{e}\cos\theta)$
  with $\hat{a}=1.524$\,AU, $\hat{e}=0.0934$.
  All residuals lie within Brahe's $2'$ precision ($\approx 0.002$\,AU).}
\label{tab:nls_residuals}
\begin{tabular}{@{}ccccc@{}}
\toprule
$i$ & $\theta_{i}$ & $r_{i}$ (AU) & $\hat{r}_{i}$ (AU) & $\hat\varepsilon_{i}$ (AU) \\
\midrule
1 &   0\textdegree & 1.382 & 1.381 & $+$0.001 \\
2 &  75\textdegree & 1.477 & 1.475 & $+$0.002 \\
3 & 150\textdegree & 1.641 & 1.643 & $-$0.002 \\
4 & 225\textdegree & 1.619 & 1.618 & $+$0.001 \\
5 & 300\textdegree & 1.442 & 1.443 & $-$0.001 \\
\midrule
\multicolumn{4}{l}{$SS_{\text{res}}=\sum\hat\varepsilon_{i}^{2}$} &
  $1.1\times10^{-5}$\,AU$^{2}$ \\
\multicolumn{4}{l}{$\text{RMSE}=\sqrt{SS_{\text{res}}/n}$} & 0.00148\,AU \\
\bottomrule
\end{tabular}
\end{table}

The RMSE of $0.00148$\,AU corresponds to an angular error of approximately
\[
  \frac{0.00148}{1.524} \times \frac{180^{\circ}}{\pi} \times 60 \approx 3.3'
\]
at Mars's mean distance.
This lies inside Brahe's observational noise floor of $2'$--$3'$, confirming
that the ellipse with two parameters absorbs all the systematic structure in
the data.

\subsection{Historical Significance of the Residual Reduction}

The eight-minute rejection of Section~\ref{sec:8min} and the NLS convergence
of this section are two sides of the same statistical coin.
Before the ellipse: $\text{RMSE}\approx4'$, four times the noise floor.
After the ellipse: $\text{RMSE}\approx2'$--$3'$, inside the noise floor.
The improvement in $SS_{\text{res}}$ from the circular to the elliptical
model is a factor of roughly $4^{2}=16$, at the cost of enlarging the
parameter space from one (the radius $a$) to two $(a,e)$.
Any model-selection criterion~---adjusted $R^{2}$, AIC, BIC, MDL~---awards
this trade-off overwhelmingly to the ellipse.

%======================================================================
\section{Maximum Likelihood and the Cram\'{e}r--Rao Bound}
\label{sec:mle}

\subsection{The Gaussian Noise Model}

Assuming independent Gaussian measurement errors,
$r_{i}= f(\theta_{i};a,e)+\varepsilon_{i}$ with $\varepsilon_{i}
\overset{iid}{\sim}\mathcal{N}(0,\sigma^{2})$, the log-likelihood is
\[
  \ell(a,e,\sigma^{2}\mid\mathbf{r})
  = -\frac{n}{2}\log(2\pi\sigma^{2})
    - \frac{1}{2\sigma^{2}}\sum_{i=1}^{n}\bigl[r_{i}-f(\theta_{i};a,e)\bigr]^{2}.
\]
Maximising over $(a,e)$ for fixed $\sigma^{2}$ is equivalent to minimising
$\sum[r_{i}-f(\theta_{i};a,e)]^{2}$, i.e.\ the NLS objective~\eqref{eq:nls}.

\textbf{Theorem (MLE = NLS).}
\textit{Under independent Gaussian observation errors, the maximum likelihood
estimator $(\hat{a}^{\text{MLE}},\hat{e}^{\text{MLE}})$ coincides with the
NLS estimator $(\hat{a},\hat{e})$.}

The proof is immediate from the log-likelihood above:
the Gaussian distributional assumption enters only through the quadratic
penalty, and the maximiser over $(a,e)$ is therefore the same as the
minimiser of the sum of squared residuals.

\subsection{Fisher Information and the Cram\'{e}r--Rao Bound}

Let $\boldsymbol\phi=(a,e)^{\top}$ denote the parameter vector.
The Fisher information matrix is
\[
  \mathcal{I}(\boldsymbol\phi) = \frac{1}{\sigma^{2}}\,\mathbf{J}^{\top}\mathbf{J},
\]
where $\mathbf{J}$ is the $n\times2$ Jacobian evaluated at the true
parameters (i.e.\ $\mathbf{J}$ has $i$-th row
$[\partial f/\partial a,\;\partial f/\partial e]$ evaluated at $\theta_{i}$
and the true $(a,e)$).
The Cram\'er--Rao bound states that for any unbiased estimator
$\tilde{\boldsymbol\phi}$,
\begin{equation}\label{eq:crb}
  \operatorname{Var}(\tilde{\boldsymbol\phi})
  \;\geq\;
  \mathcal{I}(\boldsymbol\phi)^{-1}
  = \sigma^{2}(\mathbf{J}^{\top}\mathbf{J})^{-1},
\end{equation}
where the inequality holds matrix-element-wise.
The NLS/MLE estimator achieves this bound asymptotically.

\subsection{Design Implications: Coverage and Variance Reduction}

The matrix $\mathbf{J}^{\top}\mathbf{J}$ depends on the distribution of the
observation angles $\{\theta_{i}\}$.
The partial derivatives are:
\[
  \frac{\partial f}{\partial a} = \frac{1-e^{2}}{1+e\cos\theta},
  \qquad
  \frac{\partial f}{\partial e} = -\frac{a(2e+(1+e^{2})\cos\theta)}{(1+e\cos\theta)^{2}}.
\]
Because $\partial f/\partial e$ depends on $\cos\theta$, observations near
perihelion ($\theta=0$) and aphelion ($\theta=180^{\circ}$) provide maximum
information about $e$: the planet's radial excursion from a circle is largest
near the apsides.
Observations bunched near $\theta=90^{\circ}$ or $\theta=270^{\circ}$
contribute less to eccentricity estimation because the ellipse and a circle
of the right radius agree closely near the semi-latus rectum.

Table~\ref{tab:design} summarises the design implications of
bound~\eqref{eq:crb}.

\begin{table}[h]
\centering
\renewcommand{\arraystretch}{1.25}
\caption{Design implications of the Cram\'er--Rao bound for NLS estimation
  of $(a,e)$.
  The bound favours observations spread across all orbital longitudes,
  with particular emphasis on the apsidal directions $\theta\approx0^{\circ}$
  and $\theta\approx180^{\circ}$.}
\label{tab:design}
\begin{tabular}{@{}p{4.5cm}p{4.5cm}p{4.5cm}@{}}
\toprule
\textbf{Design choice} & \textbf{Effect on $\mathbf{J}^{\top}\mathbf{J}$}
  & \textbf{Effect on variance} \\
\midrule
More observations (large $n$) &
  All eigenvalues of $\mathbf{J}^{\top}\mathbf{J}$ increase &
  Both $\operatorname{Var}(\hat a)$ and $\operatorname{Var}(\hat e)$ fall
  as $O(1/n)$ \\
Wider angular coverage (all $\theta$) &
  Off-diagonal term $\sum \partial f/\partial a \cdot \partial f/\partial e$
  cancels; matrix becomes more diagonal &
  Parameters become better separated; covariance falls \\
Apsidal observations ($\theta\approx0^{\circ}$ or $180^{\circ}$) &
  $|\partial f/\partial e|$ is maximised, increasing the $(2,2)$ element &
  $\operatorname{Var}(\hat e)$ is minimised \\
Higher instrument precision (small $\sigma^{2}$) &
  $\mathcal{I}=\sigma^{-2}\mathbf{J}^{\top}\mathbf{J}$ inflates &
  Both variances fall proportionally \\
\bottomrule
\end{tabular}
\end{table}

Brahe's twenty years of observations, spread across two orbital periods of
Mars, provided many oppositions at a wide range of true anomalies~---
exactly the design that minimises the Cram\'er--Rao bound.
Kepler's own comment on the precision of Brahe's archive~---
``Even if the positions of Mars were off by two or three minutes I would
be unable to say whether the orbit was an oval or an ellipse''~---
is a practical statement about the bound: without sub-$2'$ precision,
the eccentricity $e=0.0934$ would not be identifiable against a circular
hypothesis.

%======================================================================
\section{Kepler's Third Law: A Log-Log Regression}
\label{sec:k3}

\subsection{The Empirical Discovery}

In \emph{Harmonices Mundi} (1619), Kepler reported a relation he had
sought for over a decade: the square of a planet's orbital period is
proportional to the cube of its semi-major axis,
\begin{equation}\label{eq:k3}
  \boxed{T^{2} = k\,a^{3}}
\end{equation}
with $k$ a universal constant (the same for all planets orbiting the Sun).
Kepler found this law by scanning tables of $T^{2}$ and $a^{3}$ for the
five planets then known until the constant ratio $T^{2}/a^{3}$ became
apparent.

\subsection{Log-Log OLS}

Taking logarithms of~\eqref{eq:k3} gives
\[
  \log T = \frac{3}{2}\log a + \tfrac12 \log k.
\]
This is a bivariate linear regression
$Y_{i} = \beta X_{i} + \alpha + \varepsilon_{i}$
with $Y_{i}=\log T_{i}$, $X_{i}=\log a_{i}$, $\beta=3/2$, and
$\alpha=\frac12\log k$.
OLS gives
\[
  \hat\beta
  = \frac{S_{XY}}{S_{XX}},
  \qquad
  \hat\alpha = \bar Y - \hat\beta\bar X,
\]
where $S_{XY}=\sum(X_{i}-\bar X)(Y_{i}-\bar Y)$ and
$S_{XX}=\sum(X_{i}-\bar X)^{2}$.

\subsection{Table~3: Kepler's Third Law for the Six Known Planets}

Table~\ref{tab:harmonic} assembles the data for Mercury through Saturn
(the six planets available to Kepler and his contemporaries) in AU and
Julian years.

\begin{table}[h]
\centering
\renewcommand{\arraystretch}{1.25}
\caption{Orbital data for the six planets known in Kepler's time.
  Periods $T$ in years, semi-major axes $a$ in AU (1\,AU $=$ Earth's
  mean distance from the Sun).  The ratio $T^{2}/a^{3}$ is
  $1.000 \pm 0.001$ for all six planets.
  Log values are natural logarithms.}
\label{tab:harmonic}
\begin{tabular}{@{}lp{1.3cm}p{1.3cm}p{1.3cm}p{1.3cm}cp{1.3cm}p{1.3cm}@{}}
\toprule
Planet & $a$ (AU) & $T$ (yr) & $T^{2}$ & $a^{3}$ & $T^{2}/a^{3}$
  & $\ln a$ & $\ln T$ \\
\midrule
Mercury & 0.387 & 0.241 & 0.0581 & 0.0580 & 1.001 & $-0.949$ & $-1.423$ \\
Venus   & 0.723 & 0.615 & 0.378  & 0.378  & 1.000 & $-0.324$ & $-0.486$ \\
Earth   & 1.000 & 1.000 & 1.000  & 1.000  & 1.000 & $\phantom{-}0.000$ & $\phantom{-}0.000$ \\
Mars    & 1.524 & 1.881 & 3.538  & 3.540  & 0.999 & $\phantom{-}0.421$ & $\phantom{-}0.632$ \\
Jupiter & 5.203 & 11.86 & 140.7  & 140.8  & 0.999 & $\phantom{-}1.649$ & $\phantom{-}2.473$ \\
Saturn  & 9.537 & 29.46 & 867.5  & 867.1  & 1.001 & $\phantom{-}2.255$ & $\phantom{-}3.383$ \\
\bottomrule
\end{tabular}
\end{table}

\subsection{OLS Computation}

With $n=6$, $\bar X = 0.509$, $\bar Y = 0.763$:
\[
  S_{XX}
  = \sum_{i=1}^{6}(X_{i}-\bar X)^{2}
  = 7.431,
  \qquad
  S_{XY}
  = \sum_{i=1}^{6}(X_{i}-\bar X)(Y_{i}-\bar Y)
  = 11.147.
\]
Hence
\[
  \hat\beta = \frac{11.147}{7.431} = 1.500,
  \qquad
  \hat\alpha = 0.763 - 1.500\times0.509 = -0.001 \approx 0.
\]
The slope estimate is $\hat\beta=1.500$, consistent to four significant
figures with the theoretical value of $3/2$.
The intercept estimate $\hat\alpha\approx0$ implies $k\approx e^{0}=1.000$
in natural units of~yr$^{2}$/AU$^{3}$.

The OLS residuals $\hat\varepsilon_{i}=Y_{i}-\hat Y_{i}$ are:
Mercury: $-0.0004$; Venus: $+0.0006$; Earth: $+0.0001$;
Mars: $0.0000$; Jupiter: $+0.0004$; Saturn: $-0.0002$.
The residual sum of squares is
$SS_{\text{res}}=7.3\times10^{-7}$, giving
\[
  R^{2} = 1 - \frac{SS_{\text{res}}}{S_{YY}} = 1 - 4.4\times10^{-8} \approx 0.99998.
\]
With $n-2=4$ degrees of freedom,
the standard error of the slope is
$\text{SE}(\hat\beta) = \hat\sigma/\sqrt{S_{XX}}
= 0.000428/2.726 = 0.000157$,
giving a $t$-statistic of $\hat\beta/\text{SE}(\hat\beta)=9{,}554$.

\subsection{Identification: What Does $k$ Measure?}

The OLS regression identifies the slope $\hat\beta=3/2$ and the
intercept $\hat\alpha\approx0$, but it does not by itself identify
the physical constant $k$.
For that, Newton's mechanics is required.
Newton showed that for a central $F\propto1/r^{2}$ force of magnitude
$F=GM_{S}m/r^{2}$, a circular orbit at radius $a$ satisfies
\[
  \frac{4\pi^{2}a}{T^{2}} = \frac{GM_{S}}{a^{2}}
  \implies T^{2} = \frac{4\pi^{2}}{GM_{S}}\,a^{3}.
\]
Thus $k=4\pi^{2}/(GM_{S})$, where $G$ is Newton's gravitational constant and
$M_{S}$ the Sun's mass.
In AU--yr units, $k=1.000\,\text{yr}^{2}/\text{AU}^{3}$, which implies
$GM_{S}=4\pi^{2}\approx39.48\,\text{AU}^{3}/\text{yr}^{2}$.

This structure is analogous to the measurement-equation identification
problem in the photoelectric effect
\citep[see][Chapter 7]{Casella2002}: OLS identifies a \emph{ratio} of
physical constants ($h/e$ in the photoelectric case, $GM_{S}$ in the
orbital case), and a second measurement is needed to extract each constant
separately.

%======================================================================
\section{Newton's Inverse Problem and the Force Law}
\label{sec:newton}

\subsection{Overview of the Inference Chain}

Newton's achievement in the \emph{Principia} \citep{NewtonPrincipia1687}
was to show that Kepler's three empirical laws are theorems, not accidents.
Starting from three observed regularities~---the area law, the elliptical
orbit shape, and the harmonic law~---Newton identified the unique force
law consistent with all three simultaneously.
This is an \emph{inverse problem}: given outputs (Kepler's laws), recover
the latent mechanism (the force law).
Table~\ref{tab:newton_chain} traces the four-step logical chain.

\begin{table}[h]
\centering
\renewcommand{\arraystretch}{1.3}
\caption{Newton's four-step deductive chain connecting Kepler's empirical
  laws to the inverse-square gravitational force.
  Each step is stated in Newton's geometric language and in modern calculus.}
\label{tab:newton_chain}
\begin{tabular}{@{}p{1.0cm}p{2.7cm}p{4.9cm}p{4.9cm}@{}}
\toprule
\textbf{Step} & \textbf{Input} & \textbf{Newton's argument} & \textbf{Modern calculus} \\
\midrule
A &
  Kepler's Second Law: equal areas in equal times &
  Areas swept by the radius vector are equal in equal times.
  Newton shows (Prop.~1, Book~I) this is equivalent to zero tangential
  force component. &
  $h = r^{2}\dot\theta = \text{const}$;
  hence $\ddot r - r\dot\theta^{2} = F_{r}$ and $F_{\theta}=0$:
  the force is central. \\
\addlinespace
B &
  Kepler's Third Law: $T^{2}\propto a^{3}$ &
  For a circular orbit, centripetal acceleration is $4\pi^{2}r/T^{2}$.
  Substituting $T^{2}=ka^{3}$ and $a=r$: acceleration $\propto r^{-2}$. &
  $F = m\cdot\frac{4\pi^{2}}{k}\cdot r^{-2}$;
  identifying $k=4\pi^{2}/(GM)$ gives $F=GM m/r^{2}$. \\
\addlinespace
C &
  Force is central and $\propto r^{-2}$ &
  Prop.~11, Book~I: conic sections are produced by an inverse-square central
  force.  Conversely, an inverse-square force implies conic orbit (Prop.~13). &
  Binet equation: $u'' + u = GM/h^{2}$
  (where $u=1/r$); solution $u\propto 1+e\cos\theta$,
  i.e.\ $r=p/(1+e\cos\theta)$: an ellipse. \\
\addlinespace
D &
  Inverse-square law from solar orbits &
  Moon test: at Earth's surface, $g=9.8\,\text{m/s}^{2}$.
  At Moon's distance $r_M=60.3\,R_\oplus$, predicted centripetal
  acceleration $g/(60.3)^{2}=0.00269\,\text{m/s}^{2}$.
  Observed: $0.00272\,\text{m/s}^{2}$ (agreement 0.9\%). &
  $g_{\text{Moon}} = GM_\oplus r_M^{-2}$
  ${}= g_{\text{surf}}(R_\oplus/r_M)^{2}$;
  universality confirmed. \\
\bottomrule
\end{tabular}
\end{table}

\subsection{The Binet Equation}

The key mathematical result in Step~C is the Binet equation,
derived from Newton's second law in polar coordinates
($F_{r}=m(\ddot r - r\dot\theta^{2})$) via the substitution $u=1/r$:
\begin{equation}\label{eq:binet}
  \frac{d^{2}u}{d\theta^{2}} + u = -\frac{F(1/u)}{mh^{2}u^{2}}.
\end{equation}
For an inverse-square force $F=-GMmu^{2}$ (attractive), the right-hand side
is the constant $GM/h^{2}$, giving the linear second-order ODE
\[
  \frac{d^{2}u}{d\theta^{2}} + u = \frac{GM}{h^{2}},
\]
whose general solution is
\[
  u(\theta) = \frac{GM}{h^{2}}\bigl[1 + e\cos(\theta-\theta_{0})\bigr],
  \qquad\text{i.e.,}\quad
  r(\theta) = \frac{p}{1+e\cos\theta},\quad
  p = \frac{h^{2}}{GM}.
\]
This is equation~\eqref{eq:ellipse} with $p=a(1-e^{2})$.
Newton's genius was to run the Binet argument in reverse: \emph{given} that
orbits are ellipses (Kepler's First Law), the Binet equation implies the
right-hand side is constant, hence $F\propto u^{2}=r^{-2}$.

%======================================================================
\section{The Moon Test: Out-of-Sample Generalization}
\label{sec:moontest}

\subsection{The Numerical Argument}

The Moon test is Newton's out-of-sample validation of the inverse-square law.
Armed with Picard's (1671) improved measurement of Earth's
radius~\citep{Picard1671} ($R_\oplus = 6{,}371$\,km),
Newton computed the Moon's centripetal acceleration from its measured orbital
period and radius, and compared it to the surface gravitational acceleration.

The Moon orbits at $r_{M}=384{,}400$\,km $=60.3\,R_\oplus$ with
period $T_{M}=27.3$\,days $=2{.}36\times10^{6}$\,s.
Its centripetal acceleration is
\begin{align}
  a_{M} &= \frac{4\pi^{2}r_{M}}{T_{M}^{2}}
          = \frac{4\pi^{2}\times3{.}844\times10^{8}}{(2{.}36\times10^{6})^{2}}
          = 0{.}00272\,\text{m/s}^{2}.\label{eq:amoon}
\end{align}
Under the inverse-square law, the ratio of surface gravity to lunar centripetal
acceleration should equal the square of the distance ratio:
\[
  \frac{g}{a_{M}} \overset{?}{=} \left(\frac{r_{M}}{R_\oplus}\right)^{2}
  = (60.3)^{2} = 3{,}636.
\]
With $g=9.80\,\text{m/s}^{2}$ and $a_{M}=0.00272\,\text{m/s}^{2}$:
\[
  \left.\frac{g}{a_{M}}\right|_{\text{observed}}
  = \frac{9.80}{0.00272} = 3{,}603.
\]
The observed ratio $3{,}603$ exceeds the predicted $3{,}636$ by
$(3{,}636-3{,}603)/3{,}636=0.9\%$~---
agreement well within the measurement uncertainty in Earth's radius
and the Moon's orbital data available to Newton.

\subsection{Significance of the Moon Test}

The Moon test is the first \emph{generalization} test in the history of
physics.  Newton's model was calibrated to solar-system data (Kepler's laws);
the Moon test applied it to a completely different system (Earth--Moon) at
seven orders of magnitude closer range.
In machine-learning terms, it is an out-of-sample evaluation on a held-out
test set drawn from a different distribution.
The fact that the model passed~---with an error under $1\%$~---provided
compelling evidence for the universality of the inverse-square law.

%======================================================================
\section{AI Reconstructions of Kepler's Laws and Newton's Gravity}
\label{sec:ai}

\subsection{AI Feynman and Symbolic Regression}

Symbolic regression (SR) searches the space of mathematical expressions~---
combinations of arithmetic operations, trigonometric functions, and powers~---
for the formula that best fits a data set subject to a parsimony penalty.
\textit{AI~Feynman} \citep{Udrescu2020} extends standard SR with
physics-inspired simplifications: it exploits dimensional analysis,
decomposition into simpler sub-functions, and neural pre-smoothing of noisy
data before the symbolic search begins.

Fed a data set of $(r,\theta)$ pairs generated from planetary orbits,
AI~Feynman recovers
\[
  r = \frac{p}{1+e\cos\theta}
\]
as the most parsimonious description, landing at the ``kink'' in the
Pareto frontier (accuracy vs.\ complexity) where the ellipse achieves
a large accuracy gain at modest complexity cost.

\subsection{AI-Newton}

The \textit{AI-Newton} system of \citet{FangEtAl2023} is concept-driven:
it works with raw multi-experiment data, identifies physical concepts
(e.g.\ ``central force'', ``conserved angular momentum''), and generalises
them into symbolic laws.
Given planetary trajectory data~---without being told Kepler's laws or
Newton's~---it rediscovered:
\begin{itemize}
  \item central-force behaviour (Step~A in Table~\ref{tab:newton_chain}),
  \item the inverse-square force law (Step~B),
  \item Newton's second law and conservation of energy.
\end{itemize}
The conceptual progression~---trajectories to central force to
inverse-square law~---mirrors Newton's own chain, executed without
human supervision.

\subsection{Deep Learning: Lemos et al.\ (2023)}

\citet{Lemos2023} trained a graph neural network on approximately thirty
years of Solar System trajectory data.
The network architecture encodes the inductive bias that interactions are
pairwise and radial.
After training, symbolic regression was applied to the learned force
representation and recovered $F\propto r^{-2}$, matching Newton's law.
The network also inferred the relative masses of the planets from trajectory
curvatures~---masses are not observed directly but enter as latent variables
in the Newtonian model.

\subsection{Pareto Frontier and MDL}

A useful organising principle for comparing models of different complexity
is the \emph{Pareto frontier} in the accuracy--complexity plane.
As one moves from circular to elliptical to higher-order models, accuracy
improves but complexity grows.
The ellipse sits at the dominant kink in the frontier: beyond it,
complexity grows steeply while accuracy improves only marginally.
This kink is the hallmark of a true physical law as opposed to an
over-fitted curve.

The information-theoretic analogue is the Minimum Description Length (MDL)
principle \citep{SchmidtLipson2009}: a good model is one that compresses
the data.
The MDL cost of a model is
\[
  L(\text{Model}) = L(\text{Hypothesis}) + L(\text{Data}\mid\text{Hypothesis}),
\]
where $L(\cdot)$ is code length in nats.
The ellipse achieves a far shorter total description length than any
polynomial or Fourier series of comparable accuracy, because its
two-parameter structure encodes the physical geometry rather than
merely fitting the data.

\subsection{Neural Smoothing}

Brahe's $2'$ observational noise adds jitter to the $(r,\theta)$ data.
Modern AI Feynman \citep{Udrescu2020} handles this by first training a
neural network to interpolate the noisy observations, then
querying the smoothed neural manifold for the symbolic search.
This ``neural denoising'' step allows SR to focus on the signal rather than
the noise~---exactly as Kepler discarded the small ($<2'$) residuals of
the ellipse and focused on the systematic $8'$ excess of the circular model.

%======================================================================
\section{Kepler and Newton as Machine Learners}
\label{sec:ml}

\subsection{The Shared Structure}

Both Kepler and Newton proceeded in a pattern that is now standard in
machine learning: they identified a hypothesis class, specified features,
chose a loss function, fitted the model to data, and validated out of sample.
What was historically novel in each step reflects the limits of their era,
not any departure from the ML framework.

\subsection{Hypothesis Class and Inductive Bias}

Kepler's \emph{implicit hypothesis class} was the family of smooth closed
curves: first circles, then ovals, then conic sections.
His \emph{inductive bias}~---the preference embedded in his search strategy~---
was for geometric parsimony.
He accepted the ellipse not because he had exhausted all alternatives but
because it was the simplest curve that explained the data.

Newton's hypothesis class was more structured: the family of central force
laws $F=F(r)$, parametrised by the function $F(\cdot)$.
His inductive bias was the requirement that \emph{all three} of Kepler's
laws be jointly satisfied.
This multi-constraint identification reduced the hypothesis class to a single
function:~$F\propto r^{-2}$.

\subsection{Feature Engineering}

Kepler's key feature-engineering decision was the choice of coordinates.
Ptolemaic astronomy used geocentric Cartesian coordinates;
Copernicus shifted to heliocentric Cartesian.
Kepler's decisive step was to use heliocentric \emph{polar} coordinates
$(r,\theta)$ with the Sun at the focus.
In this representation, the ellipse has a simple two-parameter closed form
(equation~\eqref{eq:ellipse}).
In Cartesian coordinates, the same ellipse requires five parameters
($x^{2}/a^{2}+y^{2}/b^{2}=1$ plus four transformation parameters
for position and orientation).
Choosing the right representation halved the parameter space and made the
nonlinearity manageable~---a classic feature-engineering gain.

\subsection{Loss Function and Model Selection}

Kepler's effective loss function was the maximum angular residual, not the
mean squared residual.
He rejected the circular model when its maximum prediction error
($8'$) exceeded Brahe's precision ($2'$) by a factor of four.
This is analogous to a max-norm or $L_\infty$ loss, which is appropriate
when the measurement noise has a known hard upper bound.
In modern terms, Kepler performed a model-selection step based on whether
the worst-case residual was consistent with the noise model~---a robust
form of hypothesis testing.

\subsection{Generalisation and Transfer Learning}

Newton's extension of the inverse-square law from the solar system to
the Earth--Moon system is \emph{transfer learning}: a model learnt in one
domain (heliocentric planetary orbits) is applied without re-training to a
different domain (geocentric lunar orbit).
The Moon test quantified the generalisation gap and found it to be less
than $1\%$.
Newton then extended the same law to the tides, comets, and eventually to
falling bodies near Earth's surface~---a series of zero-shot applications
across domains that differ by many orders of magnitude in mass, distance,
and velocity.

\subsection{The Hypothetico-Deductive Loop}

Table~\ref{tab:hdloop} organises the full Brahe--Kepler--Newton episode
into the hypothetico-deductive (H-D) loop
that is common to both scientific inference and modern machine learning.

\begin{table}[htbp]
\centering
\small
\renewcommand{\arraystretch}{1.2}
\caption{The Brahe--Kepler--Newton episode as a hypothetico-deductive
  (H-D) loop.  Each row maps a historical step to a corresponding concept
  in modern machine learning.}
\label{tab:hdloop}
\begin{tabular}{@{}p{3.8cm}p{5.0cm}p{5.0cm}@{}}
\toprule
\textbf{H-D step} & \textbf{Historical action} & \textbf{ML analogue} \\
\midrule
Data collection &
  Brahe measures Mars positions over 20 years at $2'$ precision;
  produces the Rudolphine Tables &
  Labelled training data with known noise level; dataset
  determines statistical power \\
Hypothesis class &
  Kepler proposes conic sections (circles, ovals, ellipses) as
  candidate orbit shapes &
  Model architecture: family of parametric curves from which the
  learner selects \\
Inductive bias &
  Parsimony: accept the simplest curve that fits the data to
  within observational noise &
  Regularisation: Occam's razor, MDL, AIC/BIC, sparsity penalty \\
Parameter estimation &
  Kepler performs $\approx70$ iterative orbital adjustments;
  converges to $a=1.524$\,AU, $e=0.0934$ for Mars &
  Nonlinear least squares / Gauss--Newton; stochastic gradient
  descent for neural networks \\
Residual diagnostics &
  Kepler identifies 8' systematic error of circular model;
  fails to explain it as noise ($4\sigma$ excess) &
  Validation loss; model misspecification tests;
  residual heteroscedasticity \\
Model expansion &
  Kepler abandons circles (null hypothesis rejected),
  adopts ellipse as minimal sufficient extension &
  Architecture search; adding hidden layers or
  parameters until validation loss stops improving \\
Compression / MDL &
  Ellipse = 2 parameters; Vicarious Hypothesis = 3 parameters;
  ellipse wins on parsimony and fit &
  Minimum description length; Kolmogorov complexity;
  bits to encode model and its residuals \\
Identification &
  Newton uses Kepler's three laws simultaneously to identify
  $F(r)=F_0 r^{-2}$ uniquely &
  Structural identification in causal models;
  using multiple experimental constraints to pin down
  a latent mechanism \\
Out-of-sample test &
  Moon test: inverse-square law at $r = 60.3\,R_\oplus$;
  agreement to 0.9\% &
  Hold-out test set; domain generalisation benchmark \\
Transfer learning &
  Newton applies the same $F\propto r^{-2}$ to tides, comets,
  terrestrial gravity~---without re-training &
  Zero-shot transfer across task domains; pre-trained foundation
  model applied without fine-tuning \\
\bottomrule
\end{tabular}
\end{table}

\subsection{What ML Cannot Replicate}

Modern AI systems can rediscover Kepler's orbit law and Newton's force law
from trajectory data (Section~\ref{sec:ai}).
What they have not yet automated is the \emph{conceptual} leap that connects
the two.

Kepler needed a new \emph{physical concept}~---that the Sun exerts a
guiding force on the planets~---to motivate the hypothesis that orbits close
around the Sun at the focus of the ellipse.
Without that concept, no computable loss function points uniquely toward
the focus form~\eqref{eq:ellipse}.

Newton needed the concept of \emph{inertia} to separate force from motion:
equal areas in equal times is not intuitively a statement about force until
one has Newton's first law.
The Binet equation is a relation between the force law and the orbit shape,
but writing it down required the concept of a continuously acting force
that curves what would otherwise be straight-line motion.

Current symbolic-regression systems recover the formulas but not these
bridging concepts.
\textit{AI-Newton} \citep{FangEtAl2023} is a partial exception~---it
identifies ``central force'' as an intermediate concept~---but it does so
by hard-coding a library of physical primitives.
A fully autonomous scientific agent would need to construct those primitives
itself, which remains an open problem.

%======================================================================
\section{Comparison of Methods}
\label{sec:comparison}

Table~\ref{tab:comparison} locates Kepler's, Newton's, and modern methods
on a common set of dimensions.

\begin{table}[h]
\centering
\renewcommand{\arraystretch}{1.3}
\caption{Comparison of Kepler, Newton, Bayesian inference, symbolic
  regression (SR), and deep learning (DL) on the planet-motion problem.}
\label{tab:comparison}
\begin{tabular}{@{}p{2.7cm}p{2.7cm}p{2.7cm}p{2.7cm}p{2.7cm}@{}}
\toprule
\textbf{Feature} & \textbf{Kepler (1609)} & \textbf{Newton (1687)}
  & \textbf{Bayesian / MCMC} & \textbf{SR / DL (AI)} \\
\midrule
Optimisation &
  Iterative geometric trials ($\approx70$ for Mars) &
  Geometric deduction from axioms &
  MCMC posterior sampling &
  Gradient descent (DL); tree search (SR) \\
Error metric &
  Maximum deviation ($>8'$ threshold) &
  Logical consistency with all three Kepler laws &
  Log-likelihood; $\chi^{2}$ &
  MSE + complexity (MDL/Pareto) \\
Model search &
  Physical intuition (conic sections) &
  Inverse problem; unique solution &
  Prior over force laws &
  Tree search (SR); neural architecture (DL) \\
Noise handling &
  $8'$ threshold; accept if $<2'$ &
  Moon test residual 0.9\% &
  Bayesian posterior integrates over noise &
  Neural smoothing + SR \\
Interpretability &
  Fully geometric; transparent &
  Fully deductive; transparent &
  Credible intervals; posterior &
  Expression recovered (SR); black-box (DL) \\
Universality &
  Five-planet induction &
  Moon test; deliberate proof &
  Marginal likelihood for model selection &
  Multi-system training data \\
\bottomrule
\end{tabular}
\end{table}

%======================================================================
\section{Conclusions}

The Brahe--Kepler--Newton episode is one of the most instructive examples
of scientific inference in history.
Examined through the lens of modern statistics and machine learning,
it exhibits a coherent pipeline:

\begin{enumerate}
  \item \textbf{Data collection} (Brahe): twenty years of $2'$-precision
        Mars observations, sufficient to distinguish an ellipse
        ($e=0.0934$) from a circle ($e=0$) at the $4\sigma$ level.

  \item \textbf{Model selection by residual diagnostics} (Kepler):
        the $8'$ systematic error of the circular model exceeded the
        noise floor by a factor of four, forcing model expansion.
        The ellipse absorbed the excess with only two parameters,
        confirming it as the Pareto-optimal fit.

  \item \textbf{Nonlinear least squares} (Kepler): the ellipse parameters
        $a=1.524$\,AU and $e=0.0934$ were converged upon by approximately
        seventy iterative adjustments~---logically equivalent to
        Gauss--Newton iterations.

  \item \textbf{Log-log OLS} (Kepler): the harmonic law $T^{2}\propto a^{3}$
        is a one-slope log-log regression with $\hat\beta=1.500$,
        $\hat\alpha\approx0$, and $R^{2}=0.99998$ over six planets~---
        a cleaner fit than almost any empirical relationship in the
        physical sciences.

  \item \textbf{Inverse problem and identification} (Newton):
        Kepler's three laws jointly identify the inverse-square force law
        $F\propto r^{-2}$ uniquely.  The Binet equation formalises
        Newton's geometric deduction in modern calculus.

  \item \textbf{Out-of-sample generalization} (Newton): the Moon test
        confirmed the inverse-square law at $r=60.3\,R_\oplus$ with
        a $0.9\%$ discrepancy~---an out-of-sample validation across seven
        orders of magnitude in distance.
\end{enumerate}

Modern AI systems reproduce all six steps automatically.
\textit{AI~Feynman} recovers the ellipse from $(r,\theta)$ data.
The Lemos et al.\ network recovers $F\propto r^{-2}$ from trajectory data.
\textit{AI-Newton} reproduces the conceptual chain from trajectories to
central force to inverse-square law without prior physics knowledge.

What the AI systems have not yet fully automated is the \emph{conceptual
architecture} that makes the problem tractable: the heliocentric polar
coordinate representation, the identification of the Sun as the force
centre, and the Newtonian concept of inertia that distinguishes
``force'' from ``sustained velocity.''
Kepler and Newton were machine learners in a deep sense, but they were also
concept engineers~---and modern AI has only partially closed that gap.

\bigskip

\bibliographystyle{plainnat}
\bibliography{kepler_newton_ML}

\end{document}
